\section{哈密顿纤维化}
\stfhIndex{I157}{ha^mi^dun^xian^wei^hua}{哈密顿纤维化（Hamiltonian fibration）}
\label{sec:hamiltonian-fibration}

在弗洛尔理论中，一类特殊的辛纤维化起着重要作用。本节讨论这种称为\emph{哈密顿纤维化}的纤维化，部分采用 (GLS96) 与 (MSa94) 的论述；更详细的说明可参见这两篇文献。本节材料的取舍，完全由后文构造弗洛尔上同调中的裤子乘积，以及证明谱不变量关于该乘积的三角不等式之需要所决定。

\begin{definition}\label{def:hamiltonian-fibration}
若辛纤维化 $\pi:E\to B$ 带有一个与 $\pi$ 相容且闭的 $2$-形式
$\Omega$，则称它是\emph{哈密顿的}。
\end{definition}

下面的定理揭示了哈密顿纤维化在辛纤维化中的本质。乍看之下，“哈密顿”这一概念颇不几何；而这个定理却对辛纤维化语境中的哈密顿性质给出了非常几何的刻画，因而相当引人注目。这里的证明与 (GLS96) 和 (MSa94) 中的证明颇为不同。我们避免使用过多背景材料，只从微分学的第一原理直接出发，并分别利用例 4.2.3 与例 4.2.5 对圆周上和曲面上辛纤维化的刻画。

\begin{theorem}\label{thm:hamiltonian-fibration-characterization}
\stfhSubIndex{I149}{ha^mi^dun}{哈密顿（Hamiltonian）}{xian^wei^hua}{纤维化（fibration）}
设 $\pi:E\to B$ 是辛纤维化，$\Gamma$ 是 $E$ 上的辛联络。则下列陈述等价：
\begin{enumerate}[label=(\arabic*)]
\item 存在\emph{闭}的相容 $2$-形式 $\Omega$，使得 $\Gamma=\Gamma_\Omega$；
\item $\Gamma$ 绕 $B$ 中任意可缩环路的和乐都是哈密顿的。
\end{enumerate}
此外，若 (2) 成立，则 $E$ 上存在唯一的闭相容 $2$-形式 $\Omega_\Gamma$，它生成 $\Gamma$ 并满足归一化条件
\[
 \pi_!(\Omega_\Gamma^{k+1})=0
 \quad\text{于 }\Omega^2(B)\text{ 中},
\]
其中 $2k$ 是纤维的维数，而 $\pi_!(\Omega_\Gamma^{k+1})$ 表示沿纤维对
$\Omega_\Gamma^{k+1}$ 积分。$E$ 上任何其他生成 $\Gamma$ 的闭相容
$2$-形式 $\Omega$，都与 $\Omega_\Gamma$ 具有关系
\[
 \Omega-\Omega_\Gamma=\pi^*\beta,
\]
其中 $\beta$ 是 $B$ 上某个由 $\Omega$ 唯一确定的闭 $2$-形式。
\end{theorem}

上面的 $2$-形式 $\Omega_\Gamma\in\Omega^2(E)$ 称为 $\Gamma$ 的\emph{耦合形式}。
\stfhIndex{I087}{ou^he^xing^shi}{耦合形式（coupling form）}

在证明这个定理之前，先用一些自然产生于哈密顿动力学的例子说明这一概念。

\begin{example}[映射圆]\label{ex:mapping-circle}
\stfhIndex{I199}{ying^she^yuan}{映射圆（mapping circle）}
考虑辛流形 $(M,\omega)$ 和辛微分同胚 $\phi:M\to M$。构造\emph{映射圆}，记为 $M_\phi$：
\[
 E=M_\phi:=\{(t,x)\in[0,1]\times M\}/\mathord\sim,
\]
其中等价关系 $\sim$ 定义为
\[
 (1,\phi(x))\sim(0,x).
\]

这个等价关系还有如下更具动力系统意味的描述。考虑由离散群生成元
\[
 (t,x)\longmapsto(1+t,\phi(x))
\]
诱导的 $\mathbb R\times M$ 上的作用。这个作用是自由的，其商恰为
$M_\phi$。有自然投影
\[
 M_\phi\longrightarrow\mathbb R/\mathbb Z\cong S^1;
 \qquad [(t,x)]\longmapsto t.
\]
为看出这个纤维化的局部平凡性，考虑开覆盖
\[
 \mathbb R/\mathbb Z=U_\alpha\cup U_\beta;\qquad
 U_\alpha=[(-1/4,3/4)],\quad U_\beta=[(1/4,5/4)],
\]
以及平凡化
\[
 \Phi_\alpha:M_\phi|_{(-1/4,3/4)}
 \longrightarrow(-1/4,3/4)\times M;\qquad
 [(t,x)]\longmapsto(t,x)
\]
和
\[
 \Phi_\beta:M_\phi|_{(1/4,5/4)}
 \longrightarrow(1/4,5/4);\qquad
 [(t,x)]\longmapsto(t,\phi(x)).
\]
注意 $U_\alpha\cap U_\beta=(1/2,1)$。转移映射
$\Phi_{\alpha\beta}:(1/2,1)\times M\to(1/2,1)\times M$ 由
$\Phi_{\alpha\beta}(t,x)=(t,\phi(x))$ 给出；因为 $\phi$ 是辛的，这个映射也是辛的。

下面定义 $\pi:E=M_\phi\to S^1$ 的一个辛联络。先考虑作用
$(t,x)\mapsto(t+1,\phi(x))$ 在 $T([0,1]\times M)$ 上的线性化，它由
\[
 ((t,x),(a,\xi))\longmapsto
 ((t+1,\phi(x)),(a,d\phi_x(\xi)))
\]
给出，其中 $((t,x),(a,\xi))\in T([0,1]\times M)$。于是切丛
$TM_\phi$ 可实现为 $T([0,1]\times M)$ 关于这个作用的商。其竖直部分为
\[
 T^vM_\phi=\{((t,x),(a,\xi))\mid a=0\}/\mathord\sim,
\]
故向量 $\partial/\partial t$ 的一个提升具有形式
\[
 \frac{\partial}{\partial t}\oplus X_t(x),
\]
其中 $x\mapsto X_t$ 是 $E_t$ 上的向量场。要使这个提升与等价关系
$\sim$ 相容，它应满足
\[
 X_{t+1}(\phi(x))=d\phi_x(X_t(x)),
 \qquad\text{或等价地}\qquad
 \phi^*X_{t+1}=X_t.
\]
沿区间 $[t_0,t_1]\subset S^1$ 的和乐
$\phi:E_{t_0}\to E_{t_1}$，通过对曲线 $e:[t_0,t_1]\to E$ 求解常微分方程
\[
 \frac{de}{dt}=\widetilde X(e(t))
\]
得到，其中 $\widetilde X$ 是 $[t_0,t_1]$ 上 $\partial/\partial t$ 的水平提升。在上述平凡化中，若写
$e(t)=(t,\gamma(t))$ 以及
$\widetilde X(e(t))=(\partial/\partial t,X_t(\gamma(t)))$，则
$\gamma$ 满足 $\dot\gamma(t)=X_t(\gamma(t))$。于是，和乐是辛的，当且仅当向量场 $X_t$ 满足
$\mathcal L_{X_t}\omega=0$，即 $d(X_t\mathbin\lrcorner\omega)=0$，也即
$X_t$ 是一个（含时）辛向量场。总之，$M_\phi\to S^1$ 上的辛联络等价于选取一个满足
\[
 \phi^*X_{t+1}=X_t
\]
的辛向量场单参数族 $\{X_t\}_{t\in[0,1]}$。

注意，这样的向量场总是存在，因为可以在 $t=0,1$ 附近令 $X_t=0$，再在
$\mathbb R$ 上作周期延拓。容易看出，族 $\{X_t\}$ 给出哈密顿联络，当且仅当
$X_t$ 是哈密顿的，即对某个含时函数 $H_t$ 有 $X_t=X_{H_t}$。利用
$\phi$ 是辛的这一事实，可知 $H_t$ 满足
\[
 \phi^*(dH_{t+1})=dH_t.
\]
接着考虑 $\mathbb R\times M$ 上的\emph{闭} $2$-形式
$\Omega=\omega+d(H_t\,dt)$。按定义，它在每条纤维上的限制都是
$\omega$。在作用 $(t,x)\mapsto(t+1,\phi(x))$ 下推前 $\Omega$，得到
\[
 \phi^*\omega+\phi^*d(H_{t+1}\,dt)
 =\omega+d(H_t\,dt),
\]
即形式 $\Omega$ 下降为 $M_\phi$ 上的一个闭 $2$-形式。由于 $S^1$ 是一维的，对 $\Omega$ 的归一化要求是空泛的。因此
\[
 \Omega=\omega+d(H_t\,dt)
\]
是与该哈密顿联络相伴的唯一耦合形式。
\end{example}

反过来，有如下结论。

\begin{theorem}\label{thm:fibration-over-circle-mapping-circle}
考虑任意以 $(M,\omega)$ 为纤维的辛纤维化 $E\to S^1$。则存在辛微分同胚 $\phi$ 以及辛纤维化同构
\[
 \Phi:E\longrightarrow M_\phi.
\]
\end{theorem}

\begin{proof}
在 $E$ 上选取辛联络 $\Gamma$，并考虑平行移动映射
$\phi_{0,t}:E_0\to E_t$，其中 $E_t=\pi^{-1}(t)$。它定义了平凡化
\[
 \Psi_+:(-1/4,3/4)\times E_0\longrightarrow E|_{(-1/4,3/4)},
 \qquad \Psi_+(t,e_0)=\phi_{0,t}(e_0),
\]
以及
\[
 \Psi_-:(1/4,5/4)\times E_1,
 \qquad \Psi_-(t,e_1)=\phi_{1,t}(e_1),
\]
这里注意 $E_1=E_0$。由定义，
\[
 \Psi_-^{-1}\circ\Psi_+(t,e_0)
 =\Psi_-^{-1}(\phi_{0,t}(e_0)).
\]
注意到
\[
 \phi_{0,t}(e_0)
 =\phi_{1,t}(\phi_{t,1}\phi_{0,t}(e_0))
 =\phi_{1,t}(\phi_{0,1}(e_0)),
\]
可得
\[
 \Psi_-^{-1}\circ\Psi_+(t,e_0)=(t,\phi_{0,1}(e_0))
\]
在 $(1/4,3/4)\times E_0$ 上成立。因此，在
$E|_{(-1/4,3/4)}$ 上置 $\Phi_+=\Psi_+^{-1}$，在
$E|_{(1/4,5/4)}$ 上置 $\Phi_-=\Psi_-^{-1}$，便得到一个平凡化；它与映射圆描述中的平凡化完全相同，其中
$\phi=\phi_{0,1}$ 是联络 $\Gamma$ 在 $0\in\mathbb R/\mathbb Z$ 处的和乐。这就完成了证明。
\end{proof}

下面这个例子对弗洛尔同调很重要。

\begin{example}[映射柱]\label{ex:mapping-cylinder}
\stfhIndex{I200}{ying^she^zhu}{映射柱（mapping cylinder）}
考虑辛微分同胚同痕 $\{\phi^s\}_{s\in[0,1]}$，并在每个时刻
$s\in[0,1]$ 施行上述映射圆构造。得到辛纤维化
$E\to[0,1]\times S^1$。$E$ 上的辛联络由辛向量场的二参数族
$\{X_{(s,t)}\}$ 给出，它满足
$(\phi^s)^*X_{(s,t+1)}=X_{(s,t)}$，并且当 $s=0,1$ 时
$X_{s,t}=0$。当该联络是哈密顿联络时，它由哈密顿量的二参数族
$H_{s,t}$ 给出，后者满足
\[
 (\phi^s)^*(dH_{s,t+1})=dH_{s,t},
\]
而一个相容 $2$-形式为
\begin{equation}
 \Omega=\omega+d(H_{s,t}\,dt).
 \label{eq:mapping-cylinder-compatible-form}
\end{equation}
现在考察归一化条件。计算得
\[
 \Omega^{n+1}
 =(n+1)\omega^n\wedge d(H_{s,t}\,dt)
 =(n+1)\frac{\partial H_{s,t}}{\partial s}
   \omega^n\wedge ds\wedge dt,
\]
而其纤维积分为
\[
 \pi_!(\Omega^{n+1})
 =(n+1)\left(\int_M
   \frac{\partial H_{s,t}}{\partial s}\omega^n\right)ds\wedge dt.
\]
因此，归一化条件 $\pi_!(\Omega^{n+1})=0$ 化为
\begin{equation}
 0=\int_M\frac{\partial H_{s,t}}{\partial s}\omega^n
 =\frac{\partial}{\partial s}\left(\int_M H_{s,t}\omega^n\right).
 \label{eq:mapping-cylinder-normalization}
\end{equation}
特别地，若考虑满足归一化条件
$\int_M H_{s,t}\omega^n=C$（对所有 $(s,t)\in[0,1]^2$）的哈密顿量，其中
$C$ 是固定常数，例如 $C=0$，则\eqref{eq:mapping-cylinder-compatible-form}给出了唯一耦合形式的表达式。
\end{example}

\begin{exercise}\label{ex:mapping-cylinder-classification}
设辛纤维化 $E\to[0,1]\times S^1$ 满足
$E|_{\{0\}\times S^1}\cong\{0,1\}\times S^1$，它在
$\{0,1\}\times S^1$ 上平凡，并配备了在 $\{0,1\}\times S^1$ 上平凡的相容联络 $\Gamma$。证明，带有这样联络的纤维化同构于由
$\{\phi_{s,t}\}_{(s,t)\in[0,1]^2}$ 得到的映射柱，其中
$\phi_{0,t}=\operatorname{id}=\phi_{1,t}$，并带有例 4.2.5 中所描述的联络。
\end{exercise}

最后，我们证明下面这个著名的安布罗斯--辛格定理（由沃伦·安布罗斯（Warren Ambrose）与伊萨多·辛格（Isadore Singer）证明）的无穷维类似物。该定理描述主丛的和乐与曲率之间的关系 (AS53)：\emph{曲率生成和乐群}。由于安布罗斯与辛格的证明本质上使用了有限维性，我们改用一个简单的坎贝尔--豪斯多夫型公式，直接给出一个不同的、无坐标的证明。这个证明适用于任何结构群为纤维微分同胚群之弗雷歇李子群的纤维丛，而不只适用于哈密顿微分同胚群。

\stfhIndex{I056}{an^bu^luo^si^xin^ge^ding^li}{安布罗斯--辛格定理（Ambrose–Singer theorem）}
\begin{proposition}\label{prop:curvature-vector-hamiltonian}
设 $\Gamma$ 满足\cref{thm:hamiltonian-fibration-characterization}中的 (2)。令
$b\in B$、$v_1,v_2\in T_bB$。对 $i=1,2$，任取定义在 $b$ 的某个邻域
$U$ 上、满足 $\eta_i(b)=v_i$ 的向量场 $\eta_i$，并以 $\eta_i^\#$ 表示其水平提升。则
$[\eta_1^\#,\eta_2^\#]^v(b)$ 在纤维 $(E_b,\omega_b)$ 上定义了一个哈密顿向量场。
\end{proposition}

\begin{proof}
回顾一般流形 $M$ 上向量场李括号的基本公式：
\stfhIndex{I072}{kan^bei^er^hao^si^duo^fu^gong^shi}{坎贝尔--豪斯多夫公式（Campbell–Hausdorff formula）}
\begin{equation}
 \exp(tX)\exp(tY)\exp(-tX)\exp(-tY)
 \exp\left(\frac{t^2}{2}[X,Y]\right)=\varphi(t),
 \label{eq:campbell-hausdorff-bracket}
\end{equation}
其中 $\exp tX$ 是向量场 $X$ 的单参数子群，并且
\begin{equation}
 \operatorname{dist}(\varphi(t),\operatorname{id})\leq Ct^3
 \label{eq:campbell-hausdorff-error}
\end{equation}
对某个仅依赖于 $X,Y,M$ 的常数 $C>0$ 成立。（证明见 (Spi79)。）

给定 $v_1,v_2\in T_bB$ 及其在 $B$ 上的局部延拓 $\eta_1,\eta_2$，其中
$\eta_i(b)=v_i$。于是有曲率
\[
 \operatorname{curv}(\eta_1,\eta_2)
 =[\eta_1^\#,\eta_2^\#]-[\eta_1,\eta_2]^\#.
\]
把\eqref{eq:campbell-hausdorff-bracket}应用于 $E$ 上的
$\eta_1^\#,\eta_2^\#$，并以 $\sqrt t$ 代换 $t$，得到
\[
\begin{aligned}
 &\exp(\sqrt t\,\eta_1^\#)
  \exp(\sqrt t\,\eta_2^\#)
  \exp(-\sqrt t\,\eta_1^\#)
  \exp(-\sqrt t\,\eta_2^\#)
  \exp\left(\frac t2[\eta_1^\#,\eta_2^\#]\right)\\
 &\qquad=\varphi(\sqrt t).
\end{aligned}
\]
这里应指出，除最后一个因子外，所有因子都是水平单参数子群的
$\sqrt t$-重参数化。在等式右侧乘以
$\exp(-(t/2)[\eta_1^\#,\eta_2^\#])
 \exp((t/2)[\eta_1,\eta_2]^\#)$，得到
\[
\begin{aligned}
 &\exp(\sqrt t\,\eta_1^\#)
  \exp(\sqrt t\,\eta_2^\#)
  \exp(-\sqrt t\,\eta_1^\#)
  \exp(-\sqrt t\,\eta_2^\#)
  \exp\left(\frac t2[\eta_1,\eta_2]^\#\right)\\
 &\qquad=\varphi(\sqrt t)
  \exp\left(-\frac t2[\eta_1^\#,\eta_2^\#]\right)
  \exp\left(\frac t2[\eta_1,\eta_2]^\#\right).
\end{aligned}
\]
另一方面，有一般公式
\[
 \exp(t(X+Y))=\exp(tX)\exp(tY)(\operatorname{Id}+O(t^2)),
\]
所以
\[
\begin{aligned}
 &\varphi(\sqrt t)
  \exp\left(-\frac t2[\eta_1^\#,\eta_2^\#]\right)
  \exp\left(\frac t2[\eta_1,\eta_2]^\#\right)\\
 &\qquad=\varphi(\sqrt t)
  \exp\left(\frac t2
   \bigl(-[\eta_1^\#,\eta_2^\#]+[\eta_1,\eta_2]^\#\bigr)\right)
  \psi(t),
\end{aligned}
\]
其中 $\psi(t)$ 是微分同胚同痕，且
$\operatorname{dist}(\psi(t),\operatorname{Id})=O(t^2)$。

由于 $-[\eta_1^\#,\eta_2^\#]+[\eta_1,\eta_2]^\#$ 是竖直的，有
\[
 \pi\circ\exp\left(\frac t2
  \bigl(-[\eta_1^\#,\eta_2^\#]+[\eta_1,\eta_2]^\#\bigr)\right)
 =\operatorname{Id},
\]
故
\[
\begin{aligned}
 \pi\circ\Bigl(&\varphi(\sqrt t)
  \exp\left(\frac t2
   \bigl(-[\eta_1^\#,\eta_2^\#]+[\eta_1,\eta_2]^\#\bigr)\right)
  \psi(t)\Bigr)(b)
 =b+O(t^2).
\end{aligned}
\]
因此，由\eqref{eq:campbell-hausdorff-error}，投影曲线
\[
\begin{aligned}
 t\longmapsto\pi\Bigl(&\varphi(\sqrt t)
  \exp\left(\frac t2
   \bigl(-[\eta_1^\#,\eta_2^\#]+[\eta_1,\eta_2]^\#\bigr)\right)
  \psi(t)\Bigr)(b)=:\alpha(t)
\end{aligned}
\]
的像到点 $b\in B$ 的距离为 $O(\sqrt{t^3})$。取充分小的
$\epsilon>0$，考虑 $0\leq s\leq\epsilon$ 时的单参数曲线族
$\alpha:[0,s]\to B$。由于
$\operatorname{dist}(b,\alpha(t))\leq C\sqrt{t^3}$，可以拼接从
$\alpha(s)$ 到 $b$、长度不超过 $C\sqrt{s^3}$ 的短测地线，把它们补成从
$b$ 出发的可缩分段光滑环路。由此，这条曲线的水平提升恰为
\[
\begin{aligned}
 &\exp(\sqrt t\,\eta_1^\#)
  \exp(\sqrt t\,\eta_2^\#)
  \exp(-\sqrt t\,\eta_1^\#)
  \exp(-\sqrt t\,\eta_2^\#)
  \exp\left(\frac t2[\eta_1,\eta_2]^\#\right),
\end{aligned}
\]
其中 $0\leq t\leq s$，并相差一个 $O(\sqrt{t^3})$ 阶误差。对每个
$s$，这些环路的和乐映射 $E_b\to E_b$ 都是哈密顿的，故 $E_b$ 上相伴哈密顿同痕在 $t=0$ 处的生成向量场必为哈密顿向量场。而直接计算给出
\[
\begin{aligned}
 &\left.\frac d{dt}\right|_{t=0}
  \exp(\sqrt t\,\eta_1^\#)
  \exp(\sqrt t\,\eta_2^\#)
  \exp(-\sqrt t\,\eta_1^\#)
  \exp(-\sqrt t\,\eta_2^\#)
  \exp\left(\frac t2[\eta_1,\eta_2]^\#\right)\\
 &\qquad=\frac12\bigl([\eta_1^\#,\eta_2^\#]-[\eta_1,\eta_2]^\#\bigr).
\end{aligned}
\]
因此 $[\eta_1^\#,\eta_2^\#]-[\eta_1,\eta_2]^\#$ 是哈密顿的，这正是所要证明的。这就完成了证明。
\end{proof}

现在可以证明\cref{thm:hamiltonian-fibration-characterization}。

\begin{proof}[定理 4.2.2 的证明]
先证明 (1) 蕴含 (2)。任取 $B$ 中可缩环路 $\gamma$。则拉回
$\gamma^*E\to S^1$ 成为带闭 $2$-形式 $\pi_2^*\Omega$ 的辛纤维化，其中
$\pi_2:\gamma^*E\to E$ 是典范投影。由例 4.2.3 对 $S^1$ 上辛纤维化的分类，可设对某个辛微分同胚 $\phi$ 有
$\gamma^*E\cong M_\phi$。

令 $w=w(s,t):[0,1]\times S^1\to B$ 是 $\gamma$ 的一个收缩，满足
\[
 w|_{s=0}=\text{常值},
 \qquad w|_{s=1}(t)=\gamma(t).
\]
考虑拉回丛 $w^*E\to[0,1]\times S^1$，以及由映射
$\pi_2:w^*E\to E$ 拉回得到的形式 $\Omega_w:=\pi_2^*\Omega$。于是偶
$(w^*E,\Omega_w)$ 这个序对定义了 $[0,1]\times S^1$ 上由闭 $2$-形式
$\Omega_w$ 诱导的纤维化。现在计算 $E$ 在 $\gamma$ 上的和乐，它等于
$w^*E$ 沿 $\partial[0,1]\times S^1$ 的和乐。由例 4.2.5 及其中的习题 4.2.6，可把带相容 $2$-形式 $\Omega_w$ 的 $w^*E$ 等同于一个映射柱，其联络由显式形式
\[
 \Omega=\omega+d(H_{s,t}\,dt)
\]
给出，其中 $H_{s,t}$ 是满足
\[
 \frac{\partial}{\partial s}\int_M H_{s,t}\omega^n=0
\]
的二参数哈密顿量族。注意，由于 $w|_{s=0}$ 为常值，归一化条件使我们可以选取 $H_{0,t}\equiv0$。归一化条件蕴含
\[
 \int_M H_{s,t}\omega^n=0
\]
对所有 $s\in[0,1]$ 成立。于是沿 $\gamma$ 的和乐恰是哈密顿方程
$\dot x=X_{H_{1,t}}(x)$ 的流，因此是哈密顿的。

下面证明 (2) 蕴含 (1)。我们将显式写出唯一的耦合形式。由于在
$E_b$ 处任何与联络 $\Gamma$ 相容的 $2$-形式在纤维上的限制都是纤维辛形式，只需指定
$\Omega(\xi,\eta^\#)$ 与 $\Omega(\eta_1^\#,\eta_2^\#)$ 的值，其中
$\xi$ 切于纤维，而 $\eta^\#$、$\eta_i^\#$ 分别是 $B$ 在 $b\in B$ 处切向量
$\eta$、$\eta_i$ 的水平提升。先置
\[
 \Omega(\xi,\eta^\#)=0,
\]
以使 $\Omega$ 与联络 $\Gamma$ 相容。为指定
$\Omega(\eta_1^\#,\eta_2^\#)$，应用\cref{prop:curvature-vector-hamiltonian}并定义
\begin{equation}
 \Omega_e(\eta_1^\#,\eta_2^\#)
 :=H_{\eta_1\eta_2}(\pi(e)),
 \label{eq:coupling-form-horizontal-values}
\end{equation}
其中 $H_{\eta_1,\eta_2}$ 是 $E_{\pi(e)}$ 上逐纤维的哈密顿函数，其相伴哈密顿向量场为曲率向量
$[\eta_1^\#,\eta_2^\#]-[\eta_1,\eta_2]^\#$。换言之，
$H_{\eta_1,\eta_2}^b$ 是 $E_b$ 上满足
\begin{equation}
 dH_{\eta_1,\eta_2}^b
 =[\eta_1^\#,\eta_2^\#]^v\mathbin\lrcorner\omega_b
 \label{eq:curvature-hamiltonian}
\end{equation}
的函数。

现在把 $B$ 上的 $2$-形式在 $(v_1,v_2)\in T_bB$ 上取值，得到
\[
 \pi_!(\Omega^{n+1})(b)(v_1,v_2)
 =\int_{E_b}(v_1)^\#\wedge(v_2)^\#\mathbin\lrcorner\Omega^{n+1}.
\]
由定义，
\[
 (v_1)^\#\wedge(v_2)^\#\mathbin\lrcorner\Omega^{n+1}
 =(n+1)\Omega(v_1^\#,v_2^\#)\Omega^n.
\]
而 $\Omega(v_1^\#,v_2^\#)=H_{v_1,v_2}(b)$，故归一化条件
$\pi_!(\Omega^{n+1})=0$ 化为
\[
 \int_{E_b}H_{v_1,v_2}(b)\omega_b^n=0.
\]
当然，在哈密顿量的这个归一化之下，$H_{v_1,v_2}$ 由
$T_bB$ 中的有序对 $(v_1,v_2)$ 唯一确定，因而耦合形式 $\Omega_\Gamma$ 也唯一确定。

最后证明如此定义的形式 $\Omega=\Omega_\Gamma$ 是闭的。因为已经知道
$\Omega$ 竖直闭，只需证明
\[
 d\Omega(\xi,\eta_1^\#,\eta_2^\#)=0
\]
对竖直的 $\xi$ 和水平的 $\eta_i^\#$ 成立。直接计算得
\[
\begin{aligned}
 d\Omega(\xi,\eta_1^\#,\eta_2^\#)
 ={}&\xi[\Omega(\eta_1^\#,\eta_2^\#)]
 +\eta_1^\#[\Omega(\eta_2^\#,\xi)]
 +\eta_2^\#[\Omega(\xi,\eta_1^\#)]\\
 &-\Omega([\xi,\eta_1^\#],\eta_2^\#)
  -\Omega([\eta_1^\#,\eta_2^\#],\xi)
  -\Omega([\eta_2^\#,\xi],\eta_1^\#).
\end{aligned}
\]
容易核对
\[
\begin{gathered}
 \xi[\Omega(\eta_1^\#,\eta_2^\#)]=\xi[H_{\eta_1,\eta_2}],\\
 \eta_1^\#[\Omega(\eta_2^\#,\xi)]
 =\eta_2^\#[\Omega(\xi,\eta_1^\#)]=0,\\
 \Omega([\eta_1^\#,\eta_2^\#],\xi)
 =\Omega([\eta_1^\#,\eta_2^\#]^v,\xi)
 =dH_{\eta_1,\eta_2}(\xi).
\end{gathered}
\]
为研究其余各项，下面的引理很有用。

\begin{lemma}\label{lem:vertical-horizontal-bracket}
设 $\xi$ 是竖直向量场，即 $d\pi(\xi)=0$；考虑 $B$ 上向量场
$\eta$ 的水平提升 $\eta^\#$。则向量场 $[\xi,\eta^\#]$ 是竖直的。
\end{lemma}

\begin{proof}
注意 $\xi$ 与 $\eta^\#$ 都是可投影的，因为
\[
 d\pi(\xi)=0,
 \qquad d\pi(\eta^\#)=\eta;
\]
也就是说，它们分别与 $B$ 上的零向量场和 $\eta$ 是 $\pi$-相关的。因此
\[
 d\pi[\xi,\eta^\#]
 =[d\pi(\xi),d\pi(\eta^\#)]=[0,\eta]=0.
\]
故 $[\xi,\eta^\#]$ 是竖直的。
\end{proof}

这个引理蕴含
\[
 \Omega([\xi,\eta_1^\#],\eta_2^\#)
 =\Omega([\eta_2^\#,\xi],\eta_1^\#)=0.
\]
因此得到
\[
 d\Omega(\xi,\eta_1^\#,\eta_2^\#)
 =\xi[H_{\eta_1,\eta_2}]-dH_{\eta_1,\eta_2}(\xi)=0.
\]
这证明了 $\Omega$ 的闭性。最后一个陈述很容易核对，我们把它的证明留给读者。至此完成\cref{thm:hamiltonian-fibration-characterization}的证明。
\end{proof}

由构造可见，耦合形式 $\Omega_\Gamma$ 与联络 $\Gamma$ 的曲率密切相关。下面的总结只是\eqref{eq:curvature-hamiltonian}的直接转述。

\begin{theorem}\label{thm:coupling-form-curvature}
设 $E\to B$ 是哈密顿纤维化，$\Gamma$ 是哈密顿联络。分别以
\[
 Y_{\eta_1\eta_2}
 =(\pi^*\operatorname{curv}(\Gamma))(\eta_1^\#,\eta_2^\#),
 \qquad
 H_{\eta_1\eta_2}=\Omega_\Gamma(\eta_1^\#,\eta_2^\#)
\]
表示 $E$ 上的一个向量场和一个函数。则
\begin{equation}
 Y_{\eta_1\eta_2}\mathbin\lrcorner\Omega=dH_{\eta_1\eta_2}
 \label{eq:coupling-curvature-hamiltonian}
\end{equation}
对任意 $\eta_1,\eta_2\in TB$ 成立。
\end{theorem}

现在把上述曲率讨论专门用于乘积纤维化
$E=\Sigma\times(M,\omega)$ 的情形，其中 $\Sigma$ 是紧二维曲面，并把它与哈密顿流族联系起来。把乘积 $E=\Sigma\times(M,\omega)$ 看作辛流形丛，其结构群为
$\Symp_0(M,\omega)$，即 $\Symp(M,\omega)$ 的恒等分支。

假设 $M$ 连通，得到正合序列
\begin{equation}
 0\longrightarrow\mathbb R\longrightarrow C^\infty(M)
 \longrightarrow\operatorname{ham}(M,\omega)\longrightarrow0,
 \label{eq:functions-hamiltonian-vector-fields}
\end{equation}
其中 $\operatorname{ham}(M,\omega)$ 是 $(M,\omega)$ 上哈密顿向量场的集合。这个序列表明，哈密顿函数与哈密顿向量场之间的对应，在\emph{模掉常数函数的加法}后唯一。它是\emph{李代数}的正合序列：$\mathbb R$ 上取平凡括号，
$C^\infty(M)$ 上取泊松括号，$\operatorname{ham}(M,\omega)$ 上取李括号。在我们的约定中，第二个映射由
\[
 h\longmapsto-X_h
\]
给出。当 $M$ 闭时，这个正合序列由取均值的映射分裂：
\[
 \int_M:C^\infty(M)\longrightarrow\mathbb R;\qquad
 h\longmapsto\int_M h\,d\mu.
\]

考虑由空间 $\Sigma$ 参数化的哈密顿量族，它诱导自然正合序列
\[
 0\longrightarrow\Omega^1(\Sigma,\mathbb R)
 \longrightarrow\Omega^1(\Sigma,C^\infty(M))
 \longrightarrow\Omega^1(\Sigma,\operatorname{ham}(M,\omega))
 \longrightarrow0,
\]
从而得到同构
\[
 \Omega^1(\Sigma,C_0^\infty(M))
 \cong\Omega^1(\Sigma,\operatorname{ham}(M,\omega)),
\]
其中记 $C_0^\infty(M)=\ker\int_M$。

现在令 $K\in\Omega^1(\Sigma,C^\infty(M))$，并以 $P_K$ 表示
$\Omega^1(\Sigma,\operatorname{ham}(M,\omega))$ 中对应的 $1$-形式。于是，对每个
$\xi\in C^\infty(T\Sigma)$，$K(\xi)$ 给出 $M$ 上的函数，从而给出 $M$ 上的哈密顿向量场
$P_K(\xi)=X_{K(\xi)}$。与 $1$-形式 $K$ 相伴的一个重要量，是
$\Sigma$ 上记作 $R_K$ 的 $2$-形式；它定义为
\begin{equation}
 R_K(\xi_1,\xi_2)
 =\xi_1[K(\xi_2)]-\xi_2[K(\xi_1)]
  -\{K(\xi_2),K(\xi_1)\}
 \label{eq:connection-curvature-rk}
\end{equation}
对两个向量场 $\xi_1,\xi_2$ 成立。这里，$\xi_1[K(\xi_2)]$ 表示把
$K(\xi_2)(z,x)$ 看作 $\Sigma$ 上的函数、固定变量 $x\in M$ 后，沿向量场
$\xi_1$ 取方向导数。从这个表达式可知 $R_K$ 在 $\Sigma$ 上是张量性的。事实上，它就是与 $K$ 相伴之联络的曲率。

\begin{exercise}[见 (Ba78)]\label{ex:flat-hamiltonian-connection}
考虑平凡丛 $[0,1]^2\times(M,\omega)\to[0,1]^2$。证明联络 $P_K$ 平坦，即
$R_K=0$，当且仅当把 $P_K$ 看作依赖于 $(s,t)$ 的向量场时，它可以积分成哈密顿同痕的二参数族
\[
 \Lambda:(s,t)\longmapsto\Lambda_{(s,t)}\in\Ham(M,\omega).
\]
\end{exercise}
